\section{TYPE-013: Missing Generic on Channel — IteratorResult}
\label{sec:type-013}
\subsection*{Metadata}
\begin{tabular}{ll}
\textbf{Issue ID} & TYPE-013 \\
\textbf{Severity} & MEDIUM \\
\textbf{Category} & Type Safety \\
\textbf{Branch} & \branchref{fix/TYPE-013-channel-iterator-result-type}
\textbf{Changes} & \branchdiff{fix/TYPE-013-channel-iterator-result-type} \\
\end{tabular}

\subsection*{Executive Summary}
The \code{Channel<T>} interface at \srcfile{src/sdk.ts:35-39} uses \code{undefined as any} in \code{finish()} and \code{pull()} to work around \code{IteratorResult<T>} expecting \code{value: T} even when \code{done: true}. This breaks the generic contract and allows consumers to see \code{undefined} when they expect \code{T}.

\subsection*{Root Cause Analysis}
The \code{as any} casts hide a type mismatch:

\begin{lstlisting}[language=TypeScript]
// BEFORE — as any casts break generic contract
interface Channel<T> {
    push(v: T): void;
    finish(): void;
    pull(): Promise<IteratorResult<T>>;
}

// In finish():
resolve({ value: undefined as any, done: true });

// In pull():
return Promise.resolve({ value: undefined as any, done: true });
\end{lstlisting}

\subsection*{Fix Implementation}
The \code{IteratorResult<T, undefined>} type properly expresses the second type argument:

\begin{lstlisting}[language=TypeScript]
// AFTER — proper IteratorResult typing
interface Channel<T> {
    push(v: T): void;
    finish(): void;
    pull(): Promise<IteratorResult<T, undefined>>;
}

function createChannel<T>(): Channel<T> {
    let resolve: ((v: IteratorResult<T, undefined>) => void) | null = null;

    return {
        finish() {
            done = true;
            if (resolve) {
                resolve({ value: undefined, done: true });
                resolve = null;
            }
        },
        pull(): Promise<IteratorResult<T, undefined>> {
            if (buffer.length) {
                return Promise.resolve({ value: buffer.shift()!, done: false });
            }
            if (done) {
                return Promise.resolve({ value: undefined, done: true });
            }
            return new Promise((r) => { resolve = r; });
        },
    };
}
\end{lstlisting}

\subsection*{Affected Files}
\begin{itemize}
    \item \srcfile{src/sdk.ts} — Channel interface, createChannel()
    \item \srcfile{src/\_\_tests\_\_/type-013.test.ts}
\end{itemize}

\subsection*{Verification}
No more \code{as any} casts needed — the type correctly allows \code{undefined} when \code{done} is \code{true}.
